% year: תשפ"ה
% type: מבחן
% semester: א
% moed: א
% number: 2ב
% points: 10
% categories: func-analysis
% source: exams/מבחנים/תשפה/מועד א תשפה.pdf

\begin{question}
מצאו מינימום ומקסימום מוחלטים של הפונקציה
\[ f(x) = e^x \cdot (x-1)^2 \]
או הסבירו מדוע אינם קיימים. נמקו.
\end{question}

\begin{hint}
שימו לב ש-$f(x) \ge 0$ לכל $x$. מה קורה כאשר $x \to \infty$?
\end{hint}

\begin{solution}
$f$ מוגדרת וגזירה על כל $\R$.
\[ f'(x) = e^x(x-1)^2 + 2e^x(x-1) = e^x(x-1)(x+1). \]
נקודות קריטיות: $x = \pm 1$. $f' > 0$ על $(-\infty,-1)$, $f'<0$ על $(-1,1)$, $f'>0$ על $(1,\infty)$.
לכן $x=-1$ מקסימום מקומי ($f(-1) = 4/e$) ו-$x=1$ מינימום מקומי ($f(1)=0$).

\textbf{מינימום מוחלט:} לכל $x$ מתקיים $e^x>0$ ו-$(x-1)^2 \ge 0$, ולכן $f(x) \ge 0 = f(1)$. לכן המינימום המוחלט הוא $0$, ומתקבל ב-$x=1$.

\textbf{מקסימום מוחלט:} $\lim_{x\to\infty} e^x(x-1)^2 = \infty$ (מכפלה של שני גורמים השואפים ל-$\infty$), ולכן $f$ אינה חסומה מלעיל ואין לה מקסימום מוחלט.
(המקסימום המקומי $4/e$ ב-$x=-1$ אינו מוחלט, למשל $f(3) = 4e^3 > 4/e$.)

\textbf{תשובה:} מינימום מוחלט $0$ בנקודה $x=1$; מקסימום מוחלט אינו קיים כי $f(x)\to\infty$ כאשר $x\to\infty$.
\end{solution}
